\begin{statement}{A.5.13}{Proved}{The conductor bound when the general fibre has $q=1$ and $p_g=0$}
Let $f:Y\to\PP^1$ be a proper connected-fibre map from
a smooth compact complex threefold with
$H^1(Y,\Z)=H^2(Y,\Z)=0$. Suppose its smooth general
fibre has $h^1(\OO)=1$ and $h^2(\OO)=0$.
Then
\[
 R^1f_*\OO_Y=\OO_{\PP^1},\qquad R^2f_*\OO_Y=0,
 \qquad\Pic(Y)=\C\beta,                            \tag{A.Q1.1}
\]
where the last identification uses the exponential
coordinates of D.1.9 and its nonzero transgression
class $\beta$. Every full fibre has structure-sheaf
cohomology dimensions $(1,1,0)$.
For an irreducible reduced support $S$, every
nonconstant ambient Picard pullback to a normalization
resolution forces every reduced conductor component
to have arithmetic genus at most one. In particular
each positive-genus component is smooth elliptic,
without any lci or conductor-reducedness hypothesis.
\end{statement}
\begin{proof}
By E4, $E=R^1f_*\OO_Y$ is a line bundle. The section
induced by D.1.9's integral transgression class is
nowhere zero: its image under the natural base-change
map at each point is the nonzero coefficient class on
that full fibre. A zero of the section in $E$ would
have zero image under that map. Thus $E\simeq\OO$.
Coherent Leray now gives $H^1(Y,\OO_Y)=\C\beta$.
The exponential sequence and the two stated integral
cohomology vanishings give the Picard identification
in (A.Q1.1).

The coherent sheaf $G=R^2f_*\OO_Y$ has generic rank
zero, hence is a torsion sheaf. Riemann--Roch on $Y$
gives $\chi(\OO_Y)=\int_Yc_1c_2/24=0$, since
$c_1(Y)$ is zero in rational cohomology. The Euler
characteristic of coherent Leray is therefore
\[
 0=\chi(\OO_{\PP^1})-\chi(E)+\chi(G)
      =\operatorname{length}G.
\]
Thus $G=0$. Top-degree base change gives
$H^2(F,\OO_F)=0$ for every full Cartier fibre $F$.
The multiplication-by-parameter sequence then gives
$H^1(F,\OO_F)=E\otimes\C(c)$, of dimension one;
E4 gives $h^0(F,\OO_F)=1$.

For a full fibre $F=mS$ with irreducible reduced
support, the quotient $\OO_F\to\OO_S$ induces a
surjection on $H^2$, because its kernel has support
of dimension two and has no $H^3$.
Hence $h^2(S,\OO_S)=0$. Riemann--Roch in the Cartier
sequence on $Y$ gives $\chi(\OO_S)=0$ as in A.5.12,
so $h^1(S,\OO_S)=1$.

Put $\mathcal Q=\eta_*\OO_T/\OO_S$ for the
normalization $T$ and write $\Gamma$ for the conductor
image scheme. Nonconstant Picard pullback has nonzero
differential factoring through
$H^1(S,\OO_S)\to H^1(T,\OO_T)$.
The source has dimension one, so this map has zero
kernel. The normalization sequence (A.CD.3) gives
$H^0(\mathcal Q)=0$. Conductor duality and Serre
duality give
\[
 H^1(\Gamma,\omega_S|_\Gamma)=0,
 \qquad H^1(D,\omega_S|_D)=0                      \tag{A.Q1.2}
\]
for each reduced irreducible component $D$, where
the second map is the restriction surjection in the
proof of A.5.12. Adjunction makes $\omega_S$ ambient,
so $\deg(\omega_S|_D)=0$ by $H^2(Y,\Z)=0$.
Curve Riemann--Roch yields
$0\le h^0(D,\omega_S|_D)=1-p_a(D)$.
The normalization genus and its nonnegative delta
length then give the asserted smooth elliptic
conclusion whenever that genus is positive.
\end{proof}


% A.5.14 from part_A6_critical_curves.tex

